\chapter{Chemical Data Literacy and Public Database Foundations}

\WeekHeader{3}{Chemical Data Literacy and Public Database Foundations}{Learn what public chemical databases can and cannot prove.}{Database evidence map for the student's chemical list.}

\section{Why This Matters}
Public databases are powerful, but they are not interchangeable. PubChem, KEGG, LOTUS, FEMA, FDA/SPL, MeSH, and literature records answer different questions. Students need to know which source supports which type of claim.

\section{Database Roles}
\begin{tabularx}{\textwidth}{p{0.18\textwidth}Y Y}
\toprule
\textbf{Source} & \textbf{Useful evidence} & \textbf{Important caution}\\
\midrule
PubChem & Identifiers, synonyms, properties, annotations, hazards, classifications, assays, literature links. & Records can aggregate broad text from many sources; verify source-specific context.\\
KEGG & Compound IDs, reactions, pathways, enzymes, modules, biochemical context. & Coverage is strongest for curated biochemical systems; absence is not evidence of no role.\\
LOTUS & Natural product occurrence and taxonomic context. & Occurrence records can be sparse or uneven across organisms.\\
FEMA & Flavor-use context and sensory relevance. & Flavor presence does not prove abundance or biological role in a sample.\\
FDA/SPL & Drug product, route, dosage form, and regulatory text. & Regulatory annotation may describe a formulation context rather than the chemical alone.\\
MeSH & Pharmacologic actions and biomedical terms. & Terms may reflect literature indexing rather than direct mechanism.\\
\bottomrule
\end{tabularx}

\section{Concepts}
\begin{itemize}
  \item \textbf{Identifier ambiguity}: the same name can map to multiple structures or salts.
  \item \textbf{Synonym noise}: synonyms include trade names, old names, registry names, abbreviations, and mixture names.
  \item \textbf{Annotation source}: a database field should be interpreted through its source and section.
  \item \textbf{Ground-truth extraction}: clean values must come from traceable source fields, not guesses.
  \item \textbf{Negative space}: missing data can mean no record, no query match, no source coverage, or a real absence.
\end{itemize}

\section{Guided Lab}
Students inspect three chemicals from different contexts. For each chemical, they answer:
\begin{itemize}
  \item What are the stable identifiers?
  \item What synonyms might confuse a search?
  \item What database evidence is relevant to the project?
  \item Which fields are clean enough for grouping?
  \item Which fields are only useful as narrative evidence?
\end{itemize}

\begin{rcode}
library(uafR)
data("library_data", package = "uafR")

compounds <- c("aspirin", "caffeine", "methyl salicylate")

profile <- pubchemProfile(
  compounds = compounds,
  profile = "minimal",
  cache = TRUE,
  throttle = 0.1
)

profile$identity
profile$properties
head(profile$synonyms)
\end{rcode}

\section{Independent Extension}
Students create a source evidence map:
\begin{itemize}
  \item Rows are chemicals.
  \item Columns are evidence types: identity, properties, hazard, occurrence, pathway, reaction, sensory, biomedical, literature, assay, and regulatory.
  \item Cells contain yes, no, uncertain, or not relevant.
\end{itemize}

\section{Deliverables}
\begin{tabularx}{\textwidth}{p{0.25\textwidth}YY}
\DeliverableTableHeader
Evidence map & Shows source coverage by chemical and evidence type. & CSV or spreadsheet\\
Ambiguity note & Identifies naming or identifier problems for at least three chemicals. & Research log\\
Database interpretation paragraph & Explains which sources are most relevant to the project and why. & Report draft\\
\end{tabularx}

\section{Quality Checklist}
\begin{checklistbox}{Before Week 4}
\begin{itemize}
  \item The student can explain why a PubChem annotation, KEGG pathway, and LOTUS occurrence do not mean the same thing.
  \item The chemical intake table has stable names or identifiers whenever available.
  \item The project has a plan for handling unmatched or ambiguous chemicals.
\end{itemize}
\end{checklistbox}
